\section{Experiments}
\label{sec:main-experiments}

The experimental section is intentionally narrow.
We evaluate DGM only as a causal sequence mixer, not as a CIFAR classifier,
not as a readout adapter, and not as a general test-time adaptation wrapper.
This keeps the empirical claim aligned with the theory:
DGM should help when a sequence model needs runtime distinctions that determine
future consequences.

We use two experiments.
The first is a controlled associative-recall benchmark, where the desired
behavior is exactly to write key-value information and recover it later.
The second is language modeling from scratch, where the same DGM mixer replaces
causal attention inside a small decoder.

\subsection{Experiment 1: Multi-Query Associative Recall}
\label{sec:exp-mqar}

Multi-query associative recall (MQAR) is the cleanest benchmark for the DGM
claim \citep{arora2024zoology}.
The input contains key-value pairs and later queries; the model must return the
corresponding value.
Unlike CIFAR-style adaptation, MQAR directly tests whether a sequence mixer can
create, maintain, and retrieve runtime memory.
It also lets us compare against the same class of efficient language-model
mixers used in recent sequence-model papers.

We run DGM inside the existing PoST/Zoology MQAR framework.
The comparison is state-equalized: for each $(d_{\rm model},h)$ pair, the DGM
concept budget is chosen so that its recurrent memory budget is comparable to
the existing Mamba, RWKV, GDN, GLA/RetNet, and TTT-style runs.
Training uses sequence length $512$ with a curriculum over the number of
key-value pairs.
Evaluation reports both in-distribution recall and length extrapolation at
test lengths $512,1024,2048,4096$.
We report accuracy by slice:
\[
(\text{sequence length},\ \text{number of key-value pairs},\ \text{persist vs.
capacity}).
\]

\paragraph{Primary baselines.}
The primary comparison set is:
\begin{itemize}
    \item \textbf{Transformer}: the full-attention reference mixer.
    \item \textbf{Mamba/RWKV/GDN}: recurrent-state efficient sequence mixers.
    \item \textbf{TTT}: expressive hidden-state sequence modeling through
    test-time learning style updates.
    \item \textbf{DGM}: the proposed causal graph-memory mixer with concept
    creation, local memory repair, and edge-gated retrieval.
\end{itemize}

\paragraph{What counts as success.}
DGM does not need to beat a fully optimized Transformer on every setting.
The publication-level positive result is stronger and more specific:
\begin{enumerate}
    \item DGM matches or exceeds recurrent-state baselines at equal state
    budget on MQAR recall.
    \item DGM degrades more gracefully than fixed-state mixers as query length
    or key-value load increases.
    \item The result is reported with concept count, edge count, memory budget,
    and latency, so the gain is not hidden inside unbounded state growth.
\end{enumerate}

The implementation artifact for this experiment is the DGM-only MQAR sweep in
the PoST framework:
\path{../PoST_dev/zoology/zoology/experiments/dgm_mqar.py}.
Existing PoST MQAR runs supply the non-DGM baselines, so this script adds only
the missing DGM points rather than duplicating the whole benchmark suite.

\subsection{Experiment 2: TinyDGM-LM}
\label{sec:exp-tiny-dgm-lm}

The second experiment asks whether DGM can be trained as a language-model
mixer, rather than only solving an artificial retrieval task.
We train a tiny causal decoder from scratch in two matched variants:
\begin{itemize}
    \item \textbf{Tiny Transformer}: standard causal self-attention.
    \item \textbf{TinyDGM-LM}: the attention sublayer is replaced by a DGM
    mixer with causal left-to-right graph writes.
\end{itemize}
Both models use the same tokenizer, embedding dimension, depth, optimizer,
training tokens, and evaluation split.
The reported metrics are validation negative log likelihood, perplexity,
prefill latency, decode latency, parameter count, concept count, edge count,
and estimated runtime state size.

This experiment is not meant to be a large-scale SOTA claim.
Its job is to close the conceptual gap between MQAR and language modeling:
the same DGM operation that stores and retrieves keys in MQAR must also support
next-token training as a causal mixer.
If TinyDGM-LM is competitive with the tiny Transformer at similar parameter
count while exposing a smaller or more structured runtime state, it justifies
scaling the architecture in the PoST training stack.

The implementation artifact is
\path{experiments/train_tiny_dgm_lm.py}.
For larger runs, the same architecture is registered in PoST as
\texttt{arch=dgm}; the 180M launch script is
\path{../PoST_dev/scripts/run_dgm_180m_4b.sh}.

\subsection{Ablations Kept Out of the Main Paper}
\label{sec:exp-not-in-main}

Earlier drafts included CIFAR test-time adaptation, LoRA-style adapters,
entropy minimization, readout-only language-model memory, and synthetic
refinement diagnostics.
Those experiments are removed from the main claim.
They test different questions and made the paper look like a collection of
unrelated mechanisms.
The final empirical story is therefore:
\[
\text{DGM is a dynamic graph-memory sequence mixer.}
\]
MQAR tests whether the mixer performs controlled associative recall.
TinyDGM-LM tests whether the mixer can replace attention in a causal language
model.
